\scanpage{74}
\mspage{80}
\diagram{\begin{tikzcd}[ampersand replacement=\&,column sep=small]
X^{\unsure{\prime|1}}\arrow[r,leftrightarrow,"q"]\arrow[d,"f'_1"'] \& X'_0\arrow[r,dash]\arrow[d,"f'_0"] \& X'\arrow[r,dash]\arrow[d,"f'"] \& X\arrow[d,dash,"f"]\\
Y'_1\arrow[r,leftrightarrow,"p"']\arrow[ur,phantom,"\text{cart}" description] \& Y'_0\arrow[r]\arrow[ur,phantom,"\text{cart}" description] \& Y'\arrow[r]\arrow[ur,phantom,"\text{cart}" description] \&Y
\end{tikzcd}}

Comme $p$ est une $\mathbb D$-équiv, cela résulte bien de a).

$b)\Longrightarrow a)$. Comme la condition b) est évidemment stable par chgt de base, on est \deleted{\unsure{réduit}} ramené à prouver que \deleted{$f$ lui-$=$ est une $\mathbb D$-fibration coh. Mais c'est évident} \deleted{sous la forme a) (prop. 1) de la condition en question. Comme a) est une condition \unsure{autoduale|auto-duale} (\unsure{ou} par raisonnement direct calqué sur le précédent)} si $f$ satisfait b), et si on a
\diagram{\begin{tikzcd}[ampersand replacement=\&]
\&X\arrow[d,"f"'] \& X'\arrow[l,"q"']\arrow[d,"f'"]\\
(*)\&Y \&Y'\arrow[l,"p"]\arrow[ul,phantom,"\text{cart}" description]
\end{tikzcd}}
si $p$ est une \inserted{$\mathbb D$-}équiv., alors $q$ aussi, i.e.
\[
\underset{\substack{\unsure{\shortparallel}\\H^\bullet_{\mathbb D}(Y,f_*(M_X))}}{H^{\unsure{\bullet|\cdot|\prime}}_{\mathbb D}(X,M)}
\mathrel{\mathop{\longrightarrow}\limits^\sim}
\underset{\substack{\unsure{\shortparallel}\\H^\bullet_{\unsure{\mathbb D|\cdot}}(Y',f'_*(M_{X'}))}}{H^\bullet_{\mathbb D}(X',M)}
\]
pour tt $M$ de $A=\mathbb D(e)$. Mais \deleted{\unsure{cette flèche}} $F=f_*(M_X)$ est loc. constant et se calcule ``fibre par fibre'', de $=$ pour $F'=f'_*(M_{X'})$.

% slanted
\margin{et on a donc $F'\simeq p^*F$. On est \deleted{\unsure{conduit}} ramené à admettre que \deleted{\unsure{$f$}} $p$ étant $\mathbb D$-équiv. \unsure{implique} l'iso. $H^\bullet_{\mathbb D}(Y,F)\longrightarrow H^\bullet_{\mathbb D}(Y,p^*F)$. \unsure{C'est sûrement|C'est évidemment} \deleted{\unsure{pour}} pour $F$ constant, mais est-ce vrai pour $F$ loc. constant ? Ce n'est pas évident.}

\deleted{\unsure{s'identifie}, grâce au fait que $f,f'$ sont des $\mathbb D$-fibr. coh, à \[
\deleted{\unsure{f}}\ H^\bullet_{\mathbb D}(Z,M)\longrightarrow H^\bullet_{\mathbb D}(Z',M)
\] \unsure{où} $Z$ est la fibre de $X$ en $y$, $Z'$ celle en $y'$, et la flèche est induite par $Z'\longrightarrow Z$. Comme le \unsure{carré} est cartésien, $Z'\longrightarrow Z$ est un iso, et on gagne.}
% archive sheet number: 73
